% !PRINTED none
% Part-title (Abschnitt-opening) page: "Thirteenth Chapter. Equivalence with
% respect to Number Groups." centred, rule, then the opening of § 98 "Number
% groups in the orders." with running text below. No page number is printed in
% the running head area on this page.

\begin{center}
\gesperrt{Thirteenth Chapter.}

Equivalence with respect to Number Groups.
\end{center}

\rule{0.3\textwidth}{0.4pt}

\begin{center}
\textbf{§ 98. \gesperrt{Number groups in the orders.}}
\end{center}

In order to investigate more closely the relations of the quadratic
irrational numbers to the orders of the quadratic forms, we derive
from the numbers of the field $\Omega$ \gesperrt{number groups} according to
the following points of view.

1. We take an arbitrary natural number $Q$ and exclude from the
integers in $\Omega$ in the first place all those that are not coprime
to $Q$. From the remaining ones we also take arbitrary
quotients.

The integral and fractional numbers $\eta$ so obtained we call
\gesperrt{foreign to $Q$}. Let their totality be denoted by $O$. They
form a \gesperrt{group} in so far as multiplication and division
of two of these numbers always yield numbers of the same system.

2. We now proceed in the same way with the numbers of the order $[Q]$,
that is, we exclude all numbers of this order that are not
relatively prime to $Q$, and we also form the quotients
of any two of the numbers that remain. We denote
these numbers in $\eta'$ % [sic] German source prints in, where mit is expected. and their totality by $O'$, and call them
integral and fractional numbers of the order $[Q]$. The numbers $O'$
are all contained in $O$ and likewise form a \gesperrt{group} with respect to
multiplication and division. We call it the
\gesperrt{group of the order $[Q]$}, or briefly an \gesperrt{order group}.

Every number $\eta$ in $O$ can, by multiplication by an integral
rational number $n$ \gesperrt{coprime to $Q$}, be transformed into an integer $\omega$.

If $\eta'$ is a number in $O'$, one can make the denominator
rational and finds that $n\eta'$ is an integer of the order
$[Q]$, and consequently is congruent to a rational number with respect to
the modulus $Q$.
